The in-space assembly of kilometre-scale spacecraft represents a
critical enabling capability for future deep-space exploration,
space-based solar power, and large-aperture science missions.
This paper presents a mission planning simulation tool for the
on-orbit assembly of a modular spacecraft at the Earth-Sun L4
Lagrange point, drawing on launch vehicles from Artemis Accords
partner nations.
The optimizer formulates the assembly campaign as a
finite-horizon Markov decision process over a time-expanded
network and solves the Bellman optimality equation via forward
beam search, constraining the exponential module-set state space
to tractable size while retaining near-optimal solutions.
Eight constraints are explicitly enforced: assembly DAG
precedence, EVA crew-pair and IVA-support allocation, a 60\%
crew duty-cycle cap, scale-adaptive proximity congestion
penalties, partial work-in-progress tracking across periods,
pad turnaround intervals, first-flight cost premiums, and cargo
vehicle proximity contributions during unloading.
A multi-objective cost function combines normalized launch count,
campaign duration, and program cost with user-adjustable weights,
enabling Pareto frontier analysis across the full weight space.
A catalog of nine cargo vehicles, four crew vehicles, and three
transfer stages spanning six Artemis partner nations is supported,
along with a parametric spacecraft generator covering 14 module
types in six subsystem categories.
A Flask REST API exposes the solver to a Next.js web application
providing interactive configuration, real-time Gantt playback,
and one-click PDF report generation.
An example scenario for a 5\,km NEP spacecraft under
launch-minimizing weights
$(w_1, w_2, w_3) = (5, 1, 1)$
requires 33 launches across 22.2~months at a program cost of
\$17{,}630\,M to assemble 58~modules.
Pareto analysis shows that modest shifts toward time minimization
increase launch count only marginally, while cost-optimal
solutions incur substantial schedule duration penalties.
